% SOURCE_AUTHORITY: sga5_fr_workpass.tex, lines 4796--4833 (LF-normalized unit).
% SOURCE_SHA256: 791F4EFFC5E02832D5D77ED03518C8156D6F07E4C8238B03545DB93D883FBB28
A demonstração apoiará-se no

\begin{statement}{Lema 4.3}
Sejam $Z$ um esquema noetheriano regular conexo de dimensão $2$, de características residuais primas a $\ell$, e $X$, $Y$ divisores regulares que se intersectam transversalmente em um ponto $z$, de onde se obtém um diagrama de inclusões
\[
\begin{tikzcd}[column sep=large,row sep=large]
Z-(X\cup Y) \arrow[r,"i'"] \arrow[d,hook,"j'"'] & Z-Y \arrow[d,hook,"j"]\\
Z-X \arrow[r,"i"'] & Z\\
Y-z \arrow[r,"i_1"'] \arrow[u,hook,"m'"'] & Y \arrow[u,hook,"m"] .
\end{tikzcd}
\]
Adoptamos a hipótese de pureza:
\[
(P)\qquad \R^qj_*(\mathbb Z/\ell)=0\quad\text{para }q>1.
\]
Seja, por outro lado, $L\in\Ob D^b(Z-(X\cup Y),\Lambda)$, com cohomologia construtível e localmente constante, cuja representação correspondente de $\pi_1(Z-(X\cup Y))$ se fatore por um $\ell$-grupo. Então:
\begin{enumerate}[label=\alph*)]
\item Tem-se um isomorfismo canônico
\[
i_{1!}m'^*\R j_*i'_!L\xrightarrow{\sim} m^*\R j_*(i'_!L),
\]
e o complexo $m'^*\R j_*i'_!L$ tem cohomologia construtível e localmente constante.
\item Suponhamos que $Y$ seja conexo, e seja $\bar s$ um ponto geométrico de $Y$ situado sobre um ponto fechado $s\ne z$. Então, a flecha de restrição
\begin{equation}
\R\Gamma(Y,m^*\R j_*(i'_!L))\longrightarrow (\R j_*(i'_!L))_{\bar s}
\tag{4.3.1}
\end{equation}
é nula.
\item Sob as hipóteses de b), seja $C$ um subesquema fechado de dimensão $1$ de $Z$ tal que $Y\cap C=s$, e denotemos por $\widetilde C$ a localização estrita de $C$ em $\bar s$. Então, a flecha de restrição
\begin{equation}
\R\Gamma(Z-Y,i'_!L)
\longrightarrow
\R\Gamma(\widetilde C-\bar s,i'_!(L)|\widetilde C-\bar s)
\tag{4.3.2}
\end{equation}
é nula.
\end{enumerate}
\end{statement}
